% Reviewed translation: PDF pages 276--282 / printed pages 262--268.
% The original scan is authoritative; mathematical tokens were checked against it.
\MathBookSourcePage{276}
\subsection{$H(\operatorname{curl};\Omega)$ 中的一族协调有限元}
\label{subsec:incompressible-nedelec-elements}

本小节介绍 Nédélec~[59, 60] 提出的有限元空间.
其构造要求单元交界面上切向分量恰好连续,
因而必须使用次数为 $l$ 的不完全多项式空间, $l\geq1$.
记 $\widetilde P_l$ 为 $\mathbb R^3$ 中 $l$ 次齐次多项式空间, 并令
\begin{equation}\label{eq:incompressible-nedelec-reference-space}
S_l=\{\mathbf p\in\widetilde P_l^3;
\ \mathbf p(\mathbf x)\cdot\mathbf x=0,
\ \mathbf x=(x_1,x_2,x_3)\},
\qquad
R_l=P_{l-1}^3\oplus S_l.
\end{equation}
例如, $S_1$ 的元素必为
$\mathbf p(\mathbf x)=\boldsymbol\alpha\times\mathbf x$, 基可取
\[
\mathbf p_1=(0,-x_3,x_2),
\quad\mathbf p_2=(x_3,0,-x_1),
\quad\mathbf p_3=(-x_2,x_1,0).
\]

\MathBookSourcePage{277}
$S_2$ 中每个多项式可写成
\[
\mathbf p(\mathbf x)=\sum_{i,j=1}^3
\alpha_{ij}x_j\mathbf p_i(\mathbf x).
\]
九个多项式 $x_j\mathbf p_i$ 由关系
$\sum_{i=1}^3x_i\mathbf p_i=\mathbf x\times\mathbf x=0$ 联系,
删除其中一个后其余八个线性无关. 例如可取
\[
x_1\mathbf p_1,x_2\mathbf p_1,x_3\mathbf p_1,
x_1\mathbf p_2,x_2\mathbf p_2,x_3\mathbf p_2,
x_1\mathbf p_3,x_2\mathbf p_3
\]
为 $S_2$ 的基.

\begin{Lemma}\label{lem:incompressible-nedelec-gradient-kernel}
若 $\mathbf u\in R_l$ 且
$\operatorname{curl}\mathbf u=0$, 则
$\mathbf u=\operatorname{grad}p$, 其中 $p\in P_l$.
\end{Lemma}

\begin{Proof}
每个 $f\in\widetilde P_l$ 满足
$lf=\operatorname{grad}f\cdot\mathbf x$.
又知 $\mathbf u=\operatorname{grad}p$, $p\in P_{l+1}$.
由 \eqref{eq:incompressible-nedelec-reference-space},
$\operatorname{grad}p$ 中属于 $S_l$ 的项消失,
故 $p$ 没有 $l+1$ 次项.
\end{Proof}

\begin{Remark}\label{rem:incompressible-nedelec-arbitrary-dimension}
$R_l$ 的定义及 Lemma~\ref{lem:incompressible-nedelec-gradient-kernel}
可直接推广到任意维数 $N$.
\end{Remark}

\begin{Definition}\label{def:incompressible-nedelec-moments}
令 $\kappa\subset\mathbb R^3$ 为四面体, 其边记为 $e$, 面记为 $f$.
对某个 $s>2$ 和
$\mathbf u\in W^{1,s}(\kappa)^3$, 定义三组矩:
\[
\begin{aligned}
M_e(\mathbf u)&=
\left\{\int_e(\mathbf u\cdot\boldsymbol\tau)q\,de;
\ q\in P_{l-1}(e),\ e\subset\kappa\right\},\\
M_f(\mathbf u)&=
\left\{\int_f(\mathbf u\times\mathbf n)\cdot\mathbf q\,ds;
\ \mathbf q\in P_{l-2}^2(f),\ f\subset\kappa\right\},\\
M_\kappa(\mathbf u)&=
\left\{\int_\kappa\mathbf u\cdot\mathbf q\,dx;
\ \mathbf q\in P_{l-3}^3(\kappa)\right\},
\end{aligned}
\]
其中 $\boldsymbol\tau$ 是边 $e$ 的单位切向量.
\end{Definition}

\MathBookSourcePage{278}
\begin{Remark}\label{rem:incompressible-nedelec-moment-regularity}
Definition~\ref{def:incompressible-nedelec-moments}
必须假设 $\mathbf u$ 比 $H^1(\kappa)^3$ 略光滑,
因为当 $\mathbf u$ 仅属于 $H^1(\kappa)^3$ 时 $M_e(\mathbf u)$ 无意义.
\end{Remark}

为了构造 $H(\operatorname{curl};\Omega)$ 中的协调有限元,
需证明这些矩在 $R_l$ 上单值, 且 $M_e,M_f$
完全确定多项式在面上的切向分量.

\begin{Lemma}\label{lem:incompressible-nedelec-moment-count}
Definition~\ref{def:incompressible-nedelec-moments} 中矩的总数
等于 $R_l$ 的维数
\[
N_l=\frac12l(l+2)(l+3).
\]
\end{Lemma}

\begin{Proof}
由定义,
\[
\operatorname{card}M_e=6l,
\quad
\operatorname{card}M_f=4l(l-1),
\quad
\operatorname{card}M_\kappa=\frac12l(l-1)(l-2).
\]
三者之和为 $\frac12l(l+2)(l+3)$.
另一方面, 任意 $\widetilde P_l^3$ 多项式与
$\mathbf x$ 的点积遍历 $\widetilde P_{l+1}$,
所以条件 $\mathbf p(\mathbf x)\cdot\mathbf x=0$
给出 $\dim\widetilde P_{l+1}$ 个独立条件. 因而
\[
\dim R_l=3\dim P_l-\dim\widetilde P_{l+1}
=\frac12(l+3)(l+2)(l+1)-\frac12(l+3)(l+2)
=\frac12l(l+2)(l+3).
\]
\end{Proof}

单值性在任意四面体上不易直接证明.
先证明零矩在仿射变换下保持, 再在参考四面体上工作.
令 $F_\kappa:\widehat\kappa\to\kappa$ 为
$\mathbf x=F_\kappa(\widehat{\mathbf x})
=B_\kappa\widehat{\mathbf x}+\mathbf b_\kappa$.
标量与向量分别按
\begin{align}
\widehat\phi&=\phi\circ F_\kappa,
\label{eq:incompressible-nedelec-scalar-transform}\\
\widehat{\mathbf u}&=B_\kappa^T(\mathbf u\circ F_\kappa)
\label{eq:incompressible-nedelec-covariant-transform}
\end{align}
变换.

\MathBookSourcePage{279}
法向量和切向量满足
\begin{align}
\mathbf n\circ F_\kappa
&=\frac{(B_\kappa^{-1})^T\widehat{\mathbf n}}
{\|(B_\kappa^{-1})^T\widehat{\mathbf n}\|},
\label{eq:incompressible-nedelec-normal-transform}\\
\boldsymbol\tau\circ F_\kappa
&=\frac{B_\kappa\widehat{\boldsymbol\tau}}
{\|B_\kappa\widehat{\boldsymbol\tau}\|}.
\label{eq:incompressible-nedelec-tangent-transform}
\end{align}
采用 \eqref{eq:incompressible-nedelec-covariant-transform}
的主要原因是它在一定意义下保持旋度.
令
\begin{equation}\label{eq:incompressible-nedelec-skew-matrices}
C=(c_{ij})=(\partial u_j/\partial x_i-
\partial u_i/\partial x_j),
\qquad
\widehat C=(\widehat c_{ij})
=(\partial\widehat u_j/\partial\widehat x_i-
\partial\widehat u_i/\partial\widehat x_j).
\end{equation}
则
\begin{equation}\label{eq:incompressible-nedelec-curl-transform}
C\circ F_\kappa=(B_\kappa^{-1})^T
\widehat C B_\kappa^{-1}.
\end{equation}
所以 $\operatorname{curl}\mathbf u$ 与
$\operatorname{curl}\widehat{\mathbf u}$ 同时为零.

\begin{Lemma}\label{lem:incompressible-nedelec-space-invariance}
$R_l$ 在变换
\eqref{eq:incompressible-nedelec-covariant-transform} 下不变.
\end{Lemma}

\begin{Proof}
该变换保持每个 $P_k^3$, 故只需考虑 $\mathbf u\in S_l$.
此时
\[
\widehat{\mathbf u}(\widehat{\mathbf x})
=B_\kappa^T\mathbf u(B_\kappa\widehat{\mathbf x}+\mathbf b_\kappa)
=B_\kappa^T\mathbf u(B_\kappa\widehat{\mathbf x})
+\mathbf p(\widehat{\mathbf x}),
\]
其中 $\deg\mathbf p<l$, 且
\[
B_\kappa^T\mathbf u(B_\kappa\widehat{\mathbf x})
\cdot\widehat{\mathbf x}
=\mathbf u(B_\kappa\widehat{\mathbf x})
\cdot B_\kappa\widehat{\mathbf x}=0.
\]
故 $\widehat{\mathbf u}\in R_l$. 逆变换同理.
\end{Proof}

\begin{Lemma}\label{lem:incompressible-nedelec-moment-invariance}
Definition~\ref{def:incompressible-nedelec-moments}
中的三组矩在 $\kappa$ 上全为零, 当且仅当
$\widehat{\mathbf u}$ 的相应矩在 $\widehat\kappa$ 上全为零.
\end{Lemma}

\begin{Proof}
由 \eqref{eq:incompressible-nedelec-covariant-transform},
\[
\int_\kappa\mathbf u\cdot\mathbf q\,dx
=|\det B_\kappa|
\int_{\widehat\kappa}\widehat{\mathbf u}\cdot
B_\kappa^{-1}(\mathbf q\circ F_\kappa)\,d\widehat x,
\]
故体矩的零性等价. 又
$\mathbf u\times\mathbf n\cdot\mathbf q
=-\mathbf q\times\mathbf n\cdot\mathbf u$,
且面上的任意切向向量可写成 $\mathbf p\times\mathbf n$.
结合法向变换可知面矩的零性等价;
结合切向变换可知边矩的零性等价.
\end{Proof}

\MathBookSourcePage{280}
现在证明单值性所需的边界结论.
\begin{Lemma}\label{lem:incompressible-nedelec-face-unisolvence}
$R_l$ 中向量 $\mathbf u$ 在给定面 $f$ 上的全部矩为零,
当且仅当其在 $f$ 上的切向分量为零.
\end{Lemma}

\begin{Proof}
所有条件在仿射变换下保持, 可设 $f$ 位于平面 $x_3=0$.
切向分量为
$\mathbf u_T(x_1,x_2)=(u_1(x_1,x_2,0),u_2(x_1,x_2,0))$.
条件 $M_f=0$ 和 $M_e=0$ 分别等价于
\begin{align}
\int_f\mathbf u_T\cdot\mathbf q\,dx_1dx_2&=0
&&\forall\mathbf q\in P_{l-2}^2(f),
\label{eq:incompressible-nedelec-face-moment-zero}\\
\int_e\mathbf u_T\cdot\boldsymbol\tau q\,de&=0
&&\forall q\in P_{l-1}(e).
\label{eq:incompressible-nedelec-edge-moment-zero}
\end{align}
二维 Green 公式给出
$\int_f\operatorname{curl}\mathbf u_Tq=0$
对所有 $q\in P_{l-1}(f)$ 成立, 故
$\operatorname{curl}\mathbf u_T=0$ 于 $f$.
由 Lemma~\ref{lem:incompressible-nedelec-gradient-kernel},
$\mathbf u_T=\operatorname{grad}p$, $p\in P_l(f)$.
\eqref{eq:incompressible-nedelec-edge-moment-zero}
表明 $p$ 在 $\partial f$ 上为常数, 可取 $p=0$ 于 $\partial f$.
所以 $p=\lambda_1\lambda_2\lambda_3r$,
$r\in P_{l-3}(f)$.

\MathBookSourcePage{281}
这里 $\lambda_1,\lambda_2,\lambda_3$ 是 $f$ 的重心坐标.
把 $\mathbf q=\operatorname{grad}r$ 的相应多项式代入
\eqref{eq:incompressible-nedelec-face-moment-zero}, 得 $r=0$.
\end{Proof}

\begin{Lemma}\label{lem:incompressible-nedelec-tetrahedron-unisolvence}
若 $\mathbf u\in R_l$ 在 $\kappa$ 上的全部矩为零,
则 $\mathbf u\equiv0$.
\end{Lemma}

\begin{Proof}
Lemma~\ref{lem:incompressible-nedelec-face-unisolvence} 给出
\begin{equation}\label{eq:incompressible-nedelec-zero-tangential-trace}
\mathbf u\times\mathbf n=0\quad\text{于 }\partial\kappa,
\end{equation}
而体矩给出
\begin{equation}\label{eq:incompressible-nedelec-zero-volume-moments}
\int_\kappa\mathbf u\cdot\mathbf q\,dx=0
\qquad\forall\mathbf q\in P_{l-3}^3(\kappa).
\end{equation}
转到参考单元, Green 公式与
\eqref{eq:incompressible-nedelec-zero-tangential-trace} 给出
\[
\int_{\widehat\kappa}
\operatorname{curl}\widehat{\mathbf u}\cdot\widehat{\mathbf q},d\widehat x=0,
\quad
\operatorname{curl}\widehat{\mathbf u}\cdot\widehat{\mathbf n}=0
\text{ 于 }\partial\widehat\kappa.
\]
结合参考四面体的几何和
$\operatorname{curl}\widehat{\mathbf u}\in P_{l-1}^3$,
得到 $\operatorname{curl}\widehat{\mathbf u}=0$,
再由 \eqref{eq:incompressible-nedelec-curl-transform}
得 $\operatorname{curl}\mathbf u=0$.
Lemma~\ref{lem:incompressible-nedelec-gradient-kernel}
于是给出 $\mathbf u=\operatorname{grad}p$, $p\in P_l$.
由边界条件,
$p=\lambda_1\lambda_2\lambda_3\lambda_4r$,
$r\in P_{l-4}(\kappa)$; 再由
\eqref{eq:incompressible-nedelec-zero-volume-moments} 得 $r=0$.
\end{Proof}

\begin{Remark}\label{rem:incompressible-nedelec-full-polynomial-curl-zero}
同样的论证还表明, $P_l^3$ 中任意具有零矩的向量都满足
$\operatorname{curl}\mathbf u=0$.
先由面论证得到 $\operatorname{curl}\mathbf u\cdot\mathbf n=0$
于每个面, 再用上一个证明的体论证即可.
\end{Remark}

\MathBookSourcePage{282}
结合前两条 Lemma 得到本小节开头宣称的结论.
\begin{Theorem}\label{thm:incompressible-nedelec-unisolvence}
$R_l$ 中的向量场由四面体 $\kappa$ 上三组矩
$M_e,M_f,M_\kappa$ 唯一确定.
而且, 它在给定面 $f$ 上的切向分量只依赖于该面上定义的
$M_e$ 和 $M_f$.
\end{Theorem}

\begin{Definition}\label{def:incompressible-nedelec-interpolant}
令 $\kappa$ 为任意四面体,
$\mathbf u\in W^{1,s}(\kappa)^3$, $s>2$.
其插值 $r_\kappa\mathbf u$ 是 $R_l$ 中与 $\mathbf u$
具有相同三组矩的唯一多项式.
\end{Definition}
换言之,
\[
M_\kappa(r_\kappa\mathbf u-\mathbf u)=0,
\quad M_f(r_\kappa\mathbf u-\mathbf u)=0,
\quad M_e(r_\kappa\mathbf u-\mathbf u)=0.
\]
由空间与矩的不变性,
\begin{equation}\label{eq:incompressible-nedelec-interpolation-commuting}
B_\kappa^T[(r_\kappa\mathbf u)\circ F_\kappa]
=r_{\widehat\kappa}[B_\kappa^T(\mathbf u\circ F_\kappa)].
\end{equation}

\begin{Remark}\label{rem:incompressible-nedelec-interpolation-curl-zero}
若 $\mathbf u\in W^{1,s}(\kappa)^3$ 且
$\operatorname{curl}\mathbf u=0$, 则
$\operatorname{curl}(r_\kappa\mathbf u)=0$.
若 $\mathbf u\in P_l^3$, 则
$\operatorname{curl}(\mathbf u-r_\kappa\mathbf u)=0$.
\end{Remark}

现在设 $\Omega$ 为有界多面体,
$\mathcal T_h$ 是其四面体剖分. 对 $l\geq1$, 定义
\begin{align}
M_h&=\{\boldsymbol\mu_h\in H(\operatorname{curl};\Omega);
\ \boldsymbol\mu_h|_\kappa\in R_l
\ \forall\kappa\in\mathcal T_h\},
\label{eq:incompressible-nedelec-global-space}\\
\Phi_h&=M_h\cap H_0(\operatorname{curl};\Omega),
\label{eq:incompressible-nedelec-zero-trace-space}
\end{align}
并令
\begin{equation}\label{eq:incompressible-nedelec-global-interpolant}
(r_h\mathbf u)|_\kappa=r_\kappa\mathbf u
\qquad\forall\kappa\in\mathcal T_h.
\end{equation}

\begin{Lemma}\label{lem:incompressible-nedelec-conformity}
若 $\mathbf u\in W^{1,s}(\Omega)^3$, $s>2$, 则
$r_h\mathbf u\in M_h$.
若此外 $\mathbf u\times\mathbf n|_\Gamma=0$, 则
$r_h\mathbf u\in\Phi_h$.
\end{Lemma}
其证明直接来自
Lemma~\ref{lem:incompressible-nedelec-face-unisolvence}.
